\documentclass[11pt,a4paper]{article}
\usepackage[margin=2.5cm]{geometry}
\usepackage{amsmath,amssymb,bm}
\usepackage{xcolor}
\usepackage[colorlinks=true,linkcolor=blue,urlcolor=blue]{hyperref}
\usepackage{listings}
\usepackage{booktabs}
\usepackage{enumitem}

\lstset{
  basicstyle=\ttfamily\footnotesize,
  breaklines=true,
  frame=single,
  numbers=left,
  numberstyle=\tiny\color{gray},
  keywordstyle=\color{blue},
  commentstyle=\color{green!50!black},
  stringstyle=\color{purple},
}

\title{Julia port of the FSA-v2 model:\\
  closed-loop SMC$^2$ filter + GPU controller}
\author{Coding session writeup}
\date{2026-05-07}

\begin{document}
\maketitle

\section{The Julia FSA-v2 codebase}

\subsection{Repository layout}

The model lives at
\path{version_2_Julia/models/fsa_high_res/} and is organised as eight files
that mirror the Python source under \path{version_2/models/fsa_high_res/}.
The aggregator module \texttt{FSAHighRes.jl} \texttt{include}s them in the order
required by the inter-file using-clauses (PhiBurst before Simulation
because the latter reads its constants; CPU control before GPU control
because the latter re-uses some of the same RBF helpers; estimation last
because it depends on the framework's \texttt{EstimationModel}).
\vspace{4pt}

\noindent
\begin{tabular}{@{}lp{0.65\textwidth}@{}}
\toprule
File & Purpose \\
\midrule
\texttt{\_dynamics.jl} & SDE drift + state-dependent diffusion (Stage~G1
                         reparametrisation). Pure Julia; takes a NamedTuple
                         params for ergonomic \texttt{params.tau\_B} access. \\
\texttt{\_phi\_burst.jl} & Daily-$\Phi \to$ sub-daily morning-loaded burst
                          envelope. Reads \texttt{FSA\_STEP\_MINUTES} env
                          var to set \texttt{BINS\_PER\_DAY}. \\
\texttt{simulation.jl} & \texttt{INIT\_STATE}, \texttt{DEFAULT\_PARAMS},
                         the four observation samplers
                         (HR / sleep / stress / steps), circadian $C(t)$. \\
\texttt{\_plant.jl} & \texttt{StepwisePlant} struct. Holds truth params,
                     current $(B,F,A)$ state, RNG seed offset and a
                     \texttt{history} dict. Two advance methods:
                     \texttt{advance!} (takes daily $\Phi$, applies burst
                     envelope) and \texttt{advance\_subdaily!} (takes per-bin
                     $\Phi$ directly). \texttt{finalise} writes a
                     psim-format artefact. \\
\texttt{cpu\_control.jl} & CPU-side cost-functional builder. Produces a
                          \texttt{ControlBundle} of \texttt{(decoder,
                          cost\_fn, traj\_sample\_fn)} closures used by the
                          framework's CPU \texttt{run\_tempered\_smc\_loop}.
                          Currently unused in the production path
                          (the GPU controller superseded it) but kept for
                          debugging. Hosts the canonical \texttt{build\_rbf}. \\
\texttt{estimation.jl} & Priors (30 estimated + 5 frozen), the locally
                          guided (Pitt--Shephard) \texttt{propagate\_fn} that
                          fuses three Gaussian channels via sequential scalar
                          Kalman, the Bernoulli sleep weight, the
                          \texttt{align\_obs\_fn} that puts everything on a
                          \texttt{Float32} per-bin grid, and the
                          \texttt{build\_estimation\_model} factory. \\
\texttt{gpu\_pf.jl} & FSA-specific GPU bootstrap-PF kernel
                     (\texttt{fsa\_propagate\_segment\_kernel!}) +
                     \texttt{FSAGPUTargetBatched} struct. Calls into the
                     framework-shared \texttt{run\_segmented\_smc\_step!}
                     for resampling / Liu--West / OT rescue. \\
\texttt{gpu\_control.jl} & FSA-specific GPU cost kernel
                          (\texttt{fsa\_cost\_kernel!}) +
                          \texttt{FSAControlGPUTarget} struct +
                          \texttt{make\_log\_density\_fn} closure factory.
                          Hands a closure to the framework's
                          \texttt{run\_tempered\_smc\_gpu}. \\
\texttt{FSAHighRes.jl} & Module aggregator. \texttt{include}s the eight files
                          and re-exports the public symbols. \\
\bottomrule
\end{tabular}

\subsection{Stage-G1 reparametrisation in \texttt{\_dynamics.jl}}

The state $\bm{X}=(B,F,A)$ evolves under the It\^o SDE
\begin{align}
  \mathrm{d}B &= \kappa_B \cdot \frac{1+\varepsilon_A A}{1+\varepsilon_A A_{\rm typ}} \cdot \Phi \,\mathrm{d}t
                - \frac{B}{\tau_B}\,\mathrm{d}t + \sigma_B \sqrt{B(1-B)}\,\mathrm{d}W_B \\
  \mathrm{d}F &= \kappa_F \Phi\,\mathrm{d}t - \frac{1+\lambda_A A}{1+\lambda_A A_{\rm typ}} \cdot \frac{F}{\tau_F}\,\mathrm{d}t
                + \sigma_F \sqrt{F}\,\mathrm{d}W_F \\
  \mathrm{d}A &= \mu(B,F)\,A\,\mathrm{d}t - \eta A^3 \,\mathrm{d}t + \sigma_A \sqrt{A}\,\mathrm{d}W_A \\
  \mu(B,F) &= \mu_0 + \mu_B B - \mu_F F - \mu_{FF}(F-F_{\rm typ})^2,
\end{align}
with the operating-point constants $A_{\rm typ}=0.10$, $F_{\rm typ}=0.20$.
The reparametrisation absorbs three rank-deficient direction pairs
$(\kappa_B,\varepsilon_A)$, $(\mu_F,\mu_{FF})$, $(\tau_F,\lambda_A)$ into
strongly-identified directions and weakly-identified residuals. Truth
values use the effective parameters: $\kappa_B^{\rm eff}=0.01248$,
$\tau_F^{\rm eff}=7/1.1=6.364$, $\mu_F^{\rm eff}=0.26$, $\mu_0^{\rm eff}=0.036$.
The Julia \texttt{TRUTH\_PARAMS} NamedTuple is constructed from these
identities so it is bit-equivalent to Python's
\texttt{simulation.py:DEFAULT\_PARAMS}.

The substepped Euler--Maruyama integrator
\texttt{em\_step\_substepped} runs \texttt{n\_substeps=4} drift sub-steps
of length $\Delta t/4$, then applies one Wiener increment of variance
$\sigma(y_{\rm det})^2 \cdot \Delta t$ at the outer-bin boundary, then
performs boundary reflection ($B \mapsto -B$ if $B<0$, $B\mapsto 2-B$ if
$B>1$, $F\mapsto |F|$, $A\mapsto |A|$). This matches Python's
\texttt{imex\_step\_substepped} exactly.

\subsection{Burst envelope (\texttt{\_phi\_burst.jl})}

The controller plans daily $\Phi$ values; the plant runs the SDE on a
sub-daily grid (\texttt{BINS\_PER\_DAY}~$=$~\texttt{(60·24)/FSA\_STEP\_MINUTES}).
The expansion is a Gamma-shaped morning-loaded envelope:
\[
  \Phi(t) = \Phi_d \cdot e(k), \qquad
  e(k) \propto t_{\rm post}\,e^{-t_{\rm post}/\tau_h}
\]
in the wake window $[7\text{h}, 23\text{h}]$, normalised so
$\sum_k e(k)\cdot \Delta t_h = 24$. The integral preservation makes the
slow Banister dynamics (which respond only to the daily integral)
indistinguishable between bursty and smooth $\Phi$, but the fast
dynamics (CIR for $F$, Stuart--Landau for $A$) see the bursts. This
matters in benchmarks --- see \S\ref{sec:bug-bursty}.

\subsection{Stepwise plant (\texttt{\_plant.jl})}

\texttt{StepwisePlant} is a mutable Julia struct carrying
\texttt{state::Vector\{Float64\}}, \texttt{t\_bin::Int} (global bin
counter), the truth params and the per-channel observation history. Its
two driver methods are
\begin{itemize}[itemsep=2pt]
\item \texttt{advance!(plant, stride\_bins, Phi\_daily)} --- expects
  daily $\Phi$, expands to sub-daily via
  \texttt{expand\_daily\_phi\_to\_subdaily}, then runs the per-bin SDE.
\item \texttt{advance\_subdaily!(plant, Phi\_subdaily)} --- takes per-bin
  $\Phi$ values directly, bypasses the burst envelope. Added so the
  GPU controller's per-bin RBF schedule can be applied.
\end{itemize}
After advancing, \texttt{finalise} writes
\texttt{manifest.json}, \texttt{trajectory.npz}, per-channel
\texttt{obs/*.npz} and \texttt{exogenous/*.npz} so the closed-loop run
can be replayed by the Python psim consistency checks.

\subsection{Estimation model (\texttt{estimation.jl})}

The 30 estimated parameters are declared as a
\texttt{Vector\{Tuple\{Symbol, PriorType\}\}} with the same order as
Python's \texttt{PARAM\_PRIOR\_CONFIG} so the posterior cloud aligns
column-by-column. The locally-guided proposal
(\texttt{propagate\_fn}) computes a Gaussian prior predictive
$\mathcal{N}(\mu_{\rm pred}, P_{\rm pred})$ from the G1 drift and
state-dependent variance, then sequentially fuses the three Gaussian
observation channels (HR, stress, $\log(\mathrm{steps}+1)$) by scalar
Kalman updates gated by the \texttt{*\_present} masks. A Cholesky-3
factorisation samples $x_{\rm new}$, and the predictive log-marginal is
returned as the importance correction \texttt{pred\_lw}. The Bernoulli
sleep log-likelihood is added by \texttt{obs\_log\_weight\_fn}. This is
a line-by-line port of Python's
\texttt{estimation.py:propagate\_fn}.

\subsection{GPU particle filter (\texttt{gpu\_pf.jl})}

The model file contains only the FSA-specific propagate kernel ---
exactly the same locally-guided proposal as the CPU version, written
as a KernelAbstractions \texttt{@kernel} that runs in fp32 with one
thread per (chain, particle). The framework-shared kernels for
per-chain weight statistics, normalised-weights cumulative sum,
systematic resampling with Liu--West shrinkage, and Sinkhorn-based OT
rescue live in \texttt{julia/SMC2FC/src/Filtering/GPUSegmentedPF.jl}.
A driver routine \texttt{run\_segmented\_smc\_step!} composes them.
The model's only entry point is
\texttt{gpu\_log\_density\_batched(target, U)} returning per-chain
log-marginal-likelihood used by the outer SMC$^2$ HMC.

\subsection{GPU controller (\texttt{gpu\_control.jl})}

The control kernel \texttt{fsa\_cost\_kernel!} has one thread per
(chain $m$, common-random-number trial $t$). Per outer bin $k$:
\begin{enumerate}[itemsep=2pt]
\item Decode $\Phi(t_k)$ from the chain's $\bm{\theta}_m$ via the raw
  Gaussian RBF basis,
  $\Phi(t_k) = \Phi_{\max}\cdot \sigma\big(c_\Phi + \sum_a \theta_{m,a}\,\varphi_a(t_k)\big)$
  with $c_\Phi = \mathrm{logit}(\Phi_{\rm def}/\Phi_{\max})$.
\item PRE-step accumulate the cost integrand:
  $A_{\rm acc}\mathrel{+}= A\,\Delta t$,
  $\Phi_{\rm acc}\mathrel{+}= \Phi^2\,\Delta t$,
  $B_{\rm acc}\mathrel{+}= \max(F-F_{\max},0)^2\,\Delta t$.
\item Run the substepped EM SDE step with the trial-shared CRN noise.
\end{enumerate}
At the end of the bin loop the cost
$J_m=-A_{\rm acc} + \lambda_\Phi\Phi_{\rm acc} + \lambda_F B_{\rm acc}$
is written to \texttt{cost\_per\_thread[$i$]}. The outer wrapper
averages over trials. The model exposes a
\texttt{make\_log\_density\_fn(target)} closure that the framework's
generic \texttt{run\_tempered\_smc\_gpu}
(in \texttt{julia/SMC2FC/src/Control/GPUControlSMC.jl})
consumes. The framework provides parallel-chains HMC, a ChEES leapfrog
picker, FD-batched gradients and the tempered-SMC$^2$ outer loop;
none of those are FSA-specific.

\subsection{Bugs found while debugging the controller}
\label{sec:bugs}

Five issues turned up as the bench was iterated against the Python plot.
The first four are fixed in the Julia source; the fifth is identified but
not yet patched.

\paragraph{(1) RBF basis was row-normalised.} Python's
\texttt{smc2fc/control/rbf\_schedules.py} returns the raw Gaussian basis;
the original Julia port row-normalised it into a partition of unity.
Row-normalisation collapses any per-anchor $\theta$ variation into a
weighted average, so the decoded $\Phi(t)$ was effectively constant
regardless of $\bm\theta$ --- the reason the controller produced flat
$\Phi$ schedules. Removing the normalisation in both
\texttt{cpu\_control.jl:build\_rbf} and the GPU target's RBF block
moved the controller's optimum from $\Phi\approx 0.45$ flat to
$\Phi\approx 0.20$ flat (matching the spec text in \S10.2 of the
FSA-v2 outdated lecture notes).

\paragraph{(2) NaN guard reset state.} The original kernel had an
\texttt{if !isfinite(B)\ldots reset to (0.05, 0.30, 0.10)} guard at the
end of every bin. When fp32 rounding produced occasional small
non-finite intermediates, this silently zeroed out cost differences
between schedules. Removed.

\paragraph{(3) \texttt{$+\,\varepsilon_A$} inside cost diffusion.} The
plant's \texttt{noise\_scale\_fn} uses \texttt{sqrt(A + EPS\_A\_FROZEN)}
to keep the diffusion well-defined at $A=0$; the controller's
\texttt{diffusion\_state\_dep} uses \texttt{sqrt(max(A, 0))}. The Julia
port had carried the eps form into the cost kernel. Aligned to Python's
no-eps form for the controller path; plant kept on the eps form.

\paragraph{(4) POST-step cost accumulation.} The kernel was
accumulating $A_{\rm acc} \mathrel{+}= A\,\Delta t$ AFTER the EM step.
Python uses the PRE-step state (\texttt{y[2]}, not \texttt{y\_next[2]}),
i.e.\ a left-Riemann sum from $t=0$ to $t=T-\Delta t$. Switched.

\paragraph{(5) Burst envelope asymmetry between MPC plant and baseline
  plant.}
\label{sec:bug-bursty}
The benchmark calls
\texttt{advance!(base\_plant, ..., [1.0])} for the constant-$\Phi$
baseline (which expands $\Phi_{\rm daily}=1$ through the burst envelope
into per-bin morning-loaded peaks of $\sim$2--3) but
\texttt{advance\_subdaily!(mpc\_plant, phi\_per\_bin)} for the MPC plant
(which feeds the controller's smooth per-bin $\Phi$ directly, skipping
the burst envelope). A diagnostic Monte-Carlo with $n=32$ seeds
confirmed that under \emph{constant} smooth $\Phi=1$, the plant's mean
$\bar A = 0.1002$, but the spec quotes $\bar A = 0.0823$ for bursty
$\Phi_{\rm daily} = 1$. The 22\,\% gap is exactly the burst-envelope
asymmetry. The fix is to aggregate the controller's per-bin schedule
to daily $\Phi$ values and call \texttt{advance!} with the burst
envelope, so MPC and baseline run on the same fast-dynamics regime.

\subsection{Outstanding for next iteration}

\begin{itemize}[itemsep=2pt]
\item Apply the burst-envelope fix (5) so MPC and baseline trajectories
  are produced under the same fast-scale forcing.
\item Switch the controller's gradient from central FD ($h=10^{-4}$,
  17 chain-evals per gradient) to ForwardDiff dual-number autodiff.
  Per the bistable B3 module's empirical note, ForwardDiff with $d=8$
  out-performs Enzyme reverse-mode by $\sim$2.7$\times$ on a similar
  PF; the FSA controller has the same dimensionality so the same
  conclusion should hold.
\item Quantitative side-by-side cost comparison against Python at fixed
  $\bm\theta$, to confirm the Julia and Python cost surfaces are
  numerically equivalent after the bug fixes above.
\end{itemize}

\subsection{Wall-time and configuration}

The current end-to-end run on a single RTX 5090
(T~$=$~14 days, $h=15$ min, replan every stride, $n_{\rm SMC}^{\rm filter}=32$
inner-PF chains $\times$ $K=400$ particles, $n_{\rm SMC}^{\rm ctrl}=1024$
controller chains $\times$ $n_{\rm inner}=128$ trials, HMC
$\varepsilon=0.2$ / $L_{\rm ChEES}\in\{16,32,64,128,256\}$ /
\texttt{num\_mcmc}=10, target nats $=8$) completes in 813\,s. Python's
reference matches this configuration in $\sim$27 minutes. The Julia run
matches the spec qualitatively (rest-leaning $\Phi \in [0.10, 0.30]$,
F clearance from $0.30$ to $0.07$) but does not yet reproduce the
recovery$\to$overload pattern that yields the spec's headline
$+17\,\%$ mean-$A$ gain --- pending the burst-envelope fix and an
independent verification of the cost surface against Python at the
gradient level.

\end{document}
